%%%%%%%%%%  scan idx 53  |  folio 40  %%%%%%%%%%
Ainsi, relativement aux morphismes propres, et faisant abstraction des structures multiplicatives et des graduations, $A(X)$ est un foncteur covariant en $X$. Notons la formule classique
\begin{equation}\tag{2.3}
 f_*(x.f^*(y))=f_*(x).y .
\end{equation}

Nous utiliserons aussi la théorie de Chow des classes de Chern. Elle consiste à attacher, à tout fibré vectoriel $E$ sur $X$ (toujours supposée quasi-projective non singulière) des \emph{classes de Chern} $C^i(E)\in A^i(X)$, ($i\geq 1$), de façon à satisfaire aux conditions suivantes~; on pose $C^0(E)=1$ et
\begin{equation}\tag{2.4}
 \Ctil(E)=\left[\rang E,\sum_{i\geq 0}C^i(E)\right]\in \Atil(X)
\end{equation}
(voir chapitre I, \S 3 pour la définition de $\Atil(X)$). Les conditions à vérifier par une théorie des classes de Chern s'expriment alors ainsi :
\begin{equation}\tag{2.5}
 \Ctil(E)=\Ctil(E')+\Ctil(E'')
\end{equation}
si le fibré vectoriel $E$ est extension de $E'$ par $E''$ (si l'on pose $C(E)=\sum_{i\geq0}C^i(E)$, cette condition s'exprime $C(E)=C(E').C(E'')$). La formule (2.5) permet d'étendre $\Ctil$ en un homomorphisme additif de $\underline{K}(\xi(X))$ (cf. Chap. I, \S 2, exemple 2) dans $\Atil(X)$ :
\begin{equation}\tag{2.6}
 \Ctil:\underline{K}(\xi(X))\longrightarrow \Atil(X),
\end{equation}
et l'on veut que $\Ctil$ soit un \emph{$\lambda$-homomorphisme de $\lambda$-anneaux}~; cette dernière condition s'écrit aussi
\begin{equation}\tag{2.7}
 \Ctil(E\otimes F)=\Ctil(E)\Ctil(F)\qquad
 \Ctil(\Lambda^iE)=\lambda^i\Ctil(E),
\end{equation}
pour deux fibrés vectoriels $E$ et $F$ (bien entendu, $\Ctil$ est automatiquement compatible avec les augmentations). Soient $D$ un diviseur sur $X$ et $L(D)$ le fibré vectoriel qu'il définit~; on veut
